Question Two.
Question Two (A)
Circuit Simplification & Impedance Calculations
Given supply voltage: $v_s(t) = 16\sqrt{2}\sin(2\pi f t)\text{ V}$ with $f = 50\text{ Hz}$.
The RMS value of the voltage is $V_s = 16\text{ V}$. Using this as our reference phasor:
$$\mathbf{V}_s = 16\angle 0^\circ\text{ V}$$ 
Parallel section $Z_A$ (where $I_2$ and $I_3$ split):
$$Z_A = \frac{1 \cdot (-j2)}{1 - j2} = \frac{-j2(1 + j2)}{5} = 0.8 - j0.4\ \Omega$$ 
Parallel sub-combination $Z_B$ (far right, with $I_6$ and $I_7$):
$$Z_B = \frac{(j4)(-j2)}{j4 - j2} = \frac{8}{j2} = -j4\ \Omega$$ 
Total branch impedance for path 5 ($I_5$):
$$Z_5 = 4 + j2 + Z_B = 4 + j2 - j4 = 4 - j2\ \Omega$$ 
Total branch impedance for path 4 ($I_4$):
$$Z_4 = 2 + j6 - j2 = 2 + j4\ \Omega$$ 
Combined parallel impedance of branches 4 and 5 ($Z_C$):
$$Z_C = \frac{Z_4 \cdot Z_5}{Z_4 + Z_5} = \frac{(2 + j4)(4 - j2)}{(2 + j4) + (4 - j2)} = \frac{16 + j12}{6 + j2} = 3 + j1\ \Omega$$ 
Total Network Impedance ($Z_{\text{total}}$):
$$Z_{\text{total}} = Z_A + Z_C = (0.8 - j0.4) + (3 + j1) = 3.8 + j0.6\ \Omega = \mathbf{3.847\angle 8.97^\circ\ \Omega}$$ 
i. RMS values (in phasor representation) of all branch currents
Total Supply Current ($\mathbf{I}_1$):
$$\mathbf{I}_1 = \frac{\mathbf{V}_s}{Z_{\text{total}}} = \frac{16\angle 0^\circ}{3.847\angle 8.97^\circ} = \mathbf{4.159\angle -8.97^\circ\text{ A}}$$ 
Current $\mathbf{I}_2$ (through $1\ \Omega$):
$$\mathbf{I}_2 = \mathbf{3.722\angle -35.55^\circ\text{ A}}$$ 
Current $\mathbf{I}_3$ (through $-j2\ \Omega$):
$$\mathbf{I}_3 = \mathbf{1.861\angle 54.45^\circ\text{ A}}$$ 
Current $\mathbf{I}_4$ (through branch 4):
$$\mathbf{I}_4 = \mathbf{2.941\angle -53.97^\circ\text{ A}}$$ 
Current $\mathbf{I}_5$ (through branch 5):
$$\mathbf{I}_5 = \mathbf{2.941\angle 36.03^\circ\text{ A}}$$ 
Current $\mathbf{I}_6$ (through $+j4\ \Omega$):
$$\mathbf{I}_6 = \mathbf{2.941\angle -143.97^\circ\text{ A}}$$ 
Current $\mathbf{I}_7$ (through $-j2\ \Omega$):
$$\mathbf{I}_7 = \mathbf{5.882\angle 36.03^\circ\text{ A}}$$ 
ii. Phasor Diagram Guide
To draw the phasor diagram accurately on paper:
Draw $\mathbf{V}_s$ along the horizontal x-axis at $0^\circ$.
Plot $\mathbf{I}_1$ at a small angle of $-8.97^\circ$ (lagging below the x-axis).
Plot the remaining currents relative to the horizontal using their computed phase angles (e.g., $\mathbf{I}_3$ in the first quadrant at $+54.45^\circ$, $\mathbf{I}_6$ in the third quadrant at $-143.97^\circ$).
iii. Instantaneous expression for $I_1$ and Overall Power Factor
Peak current: $I_{1,\text{peak}} = 4.159 \times \sqrt{2} \approx 5.88\text{ A}$
Instantaneous Expression:
$$i_1(t) = \mathbf{5.88\sin(2\pi \cdot 50t - 8.97^\circ)\text{ A}}$$ 
Overall Power Factor:
$$\text{PF} = \cos(8.97^\circ) = \mathbf{0.988\text{ lagging}}$$ 
Question Two (B)
i. Impedance, Currents, and Phasor Diagram
Motor Equivalent Impedance ($Z_m$):
$$\theta_1 = \cos^{-1}(0.776) = 39.1^\circ$$ 
$$\vert{}Z_m\vert{} = \frac{V}{I_m} = \frac{230}{10} = 23\ \Omega \implies Z_m = 23\angle 39.1^\circ = \mathbf{17.85 + j14.51\ \Omega}$$ 
Circuit Currents after correction:
Active Power: $P = 230 \times 10 \times 0.776 = 1784.8\text{ W}$
New Supply Current: $I_{\text{new}} = \frac{1784.8}{230 \times 0.95} = \mathbf{8.17\text{ A}}$
Capacitor Current: $I_C = 10\sin(39.1^\circ) - 8.17\sin(18.19^\circ) = \mathbf{3.76\text{ A}}$
ii. Power Triangle Layout
Before Correction: Draw a right triangle with Real Power $P = 1784.8\text{ W}$ (base), Reactive Power $Q_1 = 1451\text{ VAR}$ (vertical side pointing up), and Apparent Power $S_1 = 2300\text{ VA}$ (hypotenuse at $39.1^\circ$).
After Correction: The base $P$ remains the same. The vertical height reduces to $Q_2 = 586.6\text{ VAR}$because the capacitor cancels out $Q_C = 864.4\text{ VAR}$ pointing downwards.
iii. Capacitor Ratings
kVAR Rating:
$$Q_C = \frac{230\text{ V} \times 3.76\text{ A}}{1000} = \mathbf{0.864\text{ kVAR}}$$ 
Equivalent Capacitance ($C$):
$$X_C = \frac{230}{3.76} = 61.17\ \Omega \implies C = \frac{1}{2\pi \cdot 50 \cdot 61.17} = \mathbf{52.03\ \mu\text{F}}$$ 
iv. Five (5) Benefits of Adding the Capacitor Bank
Reduces total current consumption from the main grid supply line (from $10\text{ A}$ to $8.17\text{ A}$).
Lowers power distribution losses ($I^2R$ heat losses) across connecting circuit cables.
Improves system voltage stability and regulation at the motor terminals.
Frees up system load capacity (kVA), allowing extra equipment to be supported safely by the main transformer grid.
Prevents utility penalties on utility billing due to low reactive energy utilization metrics.
